\documentclass[11pt]{article}
\usepackage[margin=1in]{geometry}
\usepackage{amsmath,amssymb,amsthm,hyperref}
\newcommand{\E}{\mathbb E}
\newcommand{\ph}{\operatorname{phase}}
\newcommand{\D}{\mathcal D}
\newtheorem{theorem}{Theorem}
\newtheorem{proposition}{Proposition}
\title{A common-tail dual and a separator for cubic Gaussian-phase rounding}
\author{Grothendieck certificate project}
\date{23 October 2026}
\begin{document}
\maketitle

\paragraph{Status.}
Neither Grothendieck endpoint changes. We prove an obstruction to an entire
rounding architecture, including its complex phases and convex mixtures. This
is not a lower bound on the constant. We also strengthen a necessary pool
screen by retaining the common signed Hermite tail. The latter screen has
uncertified numerical integrals; its floating values are not achievable scores.

\section*{1. Price the common tail before taking absolute values}
For a charge-one bounded response pair write its real mixed coefficients as
$c_k$, with $\sum_k|c_k|\le1$, and
\[
 K(\rho)=\sum_{k\ge0}c_k\rho^{2k+1},\qquad
 \D(c)=c_0-\sum_{k\ge1}|c_k|.
\]
Real coefficients can be obtained by mixing with the conjugate-reflected
pair; one must take the real parts before the score's absolute values.
Let $0\le\rho_i\le1$, $|s|\le1$, and suppose
\[
 \left|\sum_i a_i\rho_i^{2k+1}\right|\le1\quad(k\ge2).
 \tag{1}
\]
Absolute summability permits interchanging the finite price sum and the tail:
\[
 \left|\sum_i a_i\bigl(K(\rho_i)-c_0\rho_i-c_1\rho_i^3\bigr)\right|
 \le\sum_{k\ge2}|c_k|.
\]
Consequently
\[
 \boxed{\D(c)\le
 (1+\sum_i a_i\rho_i)c_0+
 (-s+\sum_i a_i\rho_i^3)c_1-\sum_i a_iK(\rho_i).}\tag{2}
\]
The right side is one linear functional of a response column. Hence a common
price bounds every convex mixture of columns, even though $\D$ itself is concave.
In particular, averaging individual upper bounds on $\D$ would not be valid.

A finite certificate for (1) is
\[
 \left|\sum_i a_i\rho_i^{2k+1}\right|\le1\quad(2\le k\le M),\qquad
 \sum_i|a_i|\rho_i^{2M+3}\le1.\tag{3}
\]
For $k>M$ each nonnegative power decreases, proving the remaining inequalities.
This retains cancellations between the prices on the finite front. It can be
strictly stronger than $\sum_i|a_i|\rho_i^5\le1$.
For midpoint errors $e_0,e_1,e_i$, the added allowance in (2) is exactly bounded by
\[
 |1+\sum_i a_i\rho_i|e_0+
 |-s+\sum_i a_i\rho_i^3|e_1+\sum_i|a_i|e_i.\tag{4}
\]
The results underlying (1)--(4), and finite-mixture assembly, are formalized in
\texttt{Grothe.Upper.SpectralScreen}; no Gaussian integral is assumed computed
by these Lean theorems.

\section*{2. An entire Gaussian-phase architecture is insufficient}
Let $Z$ and $G$ be independent proper standard complex Gaussians whose original
noise correlations are $\rho$ and $\rho^3$. Gaussian phase rounding of two
normalized linear forms in these grades has physical correlation
\[
 K_{A,B}(\rho)=H_{\rm phys}(A\rho+B\rho^3),\qquad |A|+|B|\le1,
\]
where
\[
 H_{\rm phys}(z)=\pi\left(\frac z4+\frac{z|z|^2}{32}
 +\frac{3z|z|^4}{256}+\frac{25z|z|^6}{4096}+\cdots\right).
\]
Indeed, Cauchy--Schwarz bounds the sum of the first- and third-grade inner
products by one. This also permits multiple independent Gaussians in each
grade. The identity is the Gaussian phase identity proved elsewhere in the
repository. The physical kernel is not the holomorphic odd series evaluated
at $A\rho+B\rho^3$ when these coefficients are complex.

\begin{theorem}[Uniform architecture separator]
For every $|A|+|B|\le1$, the mixed coefficients of $K_{A,B}$ satisfy
\[
 \Re(c_0-\tfrac9{25}c_1-c_2-c_3)<0.71186.
 \tag{5}
\]
Every probability mixture, and its coefficient-$\ell^1$ closure, therefore has
\[
 \boxed{\D(c)\le0.71186<
 \frac{355973576656604491}{500000000000000000}.}\tag{6}
\]
The rightmost fraction is the incumbent full reciprocal score, under its
documented analytic inputs. Thus this architecture cannot improve that score.
\end{theorem}

\begin{proof}
By charge-one equivariance factor out the common phase of $A$, and put
$A=a\ge0$, $B=rz$, $r\ge0$, $a+r\le1$, $|z|=1$.
When $A=0$, choose any common phase and use the same representation; when $r=0$
the choice of $z$ is immaterial. Expanding the physical kernel gives
\begin{align*}
 c_0/\pi&=a/4,\\
 c_1/\pi&=B/4+a^3/32,\\
 c_2/\pi&=a^2(2B+\bar B)/32+3a^5/256,\\
 c_3/\pi&=a(2|B|^2+B^2)/32
 +(9a^4B+6a^4\bar B)/256+25a^7/4096.
\end{align*}
For $e=9/25$ write
\[
 (c_0-ec_1-c_2-c_3)/\pi=F(z)=U-vz-w\bar z-tz^2,
\]
with real quantities
\begin{align*}
 U&=a/4-ea^3/32-3a^5/256-25a^7/4096-ar^2/16,\\
 v&=r(e/4+a^2/16+9a^4/256),\qquad
 w=r(a^2/32+3a^4/128),\qquad t=ar^2/32.
\end{align*}
Set $C_1=-U(v+w)+vt$, $C_2=vw-Ut$, $C_3=wt\ge0$.
The exact certificate proves
\[
 C_1-4C_2+9C_3\le0,\qquad C_1+4C_2+9C_3\le0.\tag{7}
\]
Here is a finite reproducible positivity proof for the whole triangle, including
its boundary. Divide $-(C_1\pm4C_2+9C_3)$ by $ar$ as a polynomial, substitute
$r=(1-a)b$, and express it in the tensor Bernstein basis of degrees $(10,3)$
on $[0,1]^2$. Every one of its 44 coefficients is at least $9/640$ for each
sign. Multiplying back by $ar\ge0$ proves (7); no division at a boundary point
is needed. The exported matrices are checked by exact reconstruction, and the
Lean proof expands the positive polynomials directly with \texttt{ring} and
\texttt{positivity}, without \texttt{native\_decide}.

If $x=\Re z$, direct expansion gives
\[
 |F(-1)|^2-|F(z)|^2
 =2(1+x)\{-C_1+2C_2(1-x)-C_3(2x-1)^2\}\ge0,
\]
using $-1\le x\le1$ and $C_1+4|C_2|+9C_3\le0$.
Also $F(-1)=U+v+w-t\ge0$. It increases in $r\in[0,1-a]$:
its derivative is
\[
 e/4+3a^2/32+15a^4/256-3ar/16
 \ge e/4-3/64>0.
\]
It remains to bound the univariate polynomial $P(a)=F(-1)|_{r=1-a}$.
At $a_0=29/32$ its tangent remainder is
\[
 P(a_0)+P'(a_0)(a-a_0)-P(a)=(a-a_0)^2 Q(a).
\]
The degree-five Bernstein coefficients of $Q$ are
\[
 \frac{285761973783}{1717986918400},\quad
 \frac{1920796785347}{8589934592000},\quad
 \frac{2487532348179}{8589934592000},\quad
 \frac{3202332295011}{8589934592000},\quad
 \frac{4157517803443}{8589934592000},\quad
 \frac{220673771643}{343597383680}.
\]
Each exceeds $4/25$. Completing the square gives the exact rational comparison
\[
 P(a)\le P(a_0)+\frac{25}{16}P'(a_0)^2
 <\frac{22659}{100000}.
\]
Therefore, including an arbitrary factored common phase,
\[
 \Re(c_0-ec_1-c_2-c_3)\le\pi|F(z)|
 <\frac{355}{113}\frac{22659}{100000}<0.71186.
\]
For real mixed coefficients,
$\D(c)\le c_0-ec_1-c_2-c_3$. This inequality and (5) use one linear
functional, so they apply to convex mixtures. Continuity gives the non-strict
bound in the $\ell^1$ closure. This proves (6).
\end{proof}

\paragraph{Formal boundary.}
The algebraic triangle certificate, trigonometric modulus comparison, scalar
cap, and the full-score implication from the displayed low-coefficient identity
are formalized in \texttt{Grothe.Upper.CubicPhaseScreen}. The general Gaussian
phase identity already exists in Lean. Its composition, extraction of these four
coefficients, normalization of arbitrary complex $A,B$, and the closure argument
are supplied in this written application; they have not been wired into one
end-to-end Lean theorem about the Gaussian family. No new hypothesis is added
to either Grothendieck endpoint. The separator is not a universal witness cap.

\section*{3. Exact angular averaging of the shared projection state}
For the symmetric cubic $q_m$ in the preceding replica note, set
$d=m(m+1)/2$ and $D=m(m+1)$. Its one-point state can be written using independent
$T,S\sim\operatorname{Exp}(1)$, $R\sim\Gamma(m,1)$, $V\sim\Gamma(d-1,1)$:
\[
 A_X=(R-m)/\sqrt m,\quad A_Y=(S+V-d)/\sqrt d,\quad |q_m|=R\sqrt S/\sqrt D.
\]
For an image coefficient vector $c$ define
\[
 b=c_0+c_1(T-2)/\sqrt2+c_2 A_X+c_3 A_Y,
 \qquad a=\sqrt T|b|,\quad v=|c_4|R\sqrt S/\sqrt D.
\]
Let $u(a,b)=\E_\theta\ph(a+be^{i\theta})$ for nonnegative $a,b$; this is real.
The exact relative-phase average gives the visible projection integrand
$\sqrt T\ph(b)u(a,v)$ and the remote one
$\ph(c_4)(R\sqrt S/\sqrt D)u(v,a)$. The visible factors for the other three
projections are $(T-2)/\sqrt2,A_X,A_Y$. In particular, the $S$ appearing in
$A_Y$ is the same $S$ as in $|q_m|$.

For $a<b$, $t=a/b$ gives
\[
 u(a,b)=\frac t2\,{}_2F_1(1/2,1/2;2;t^2).
\]
For $a\ge b$ and $a>0$, put $r=b/a$ and $\chi=((1-r)/(1+r))^2$. In the
parameter convention for complete elliptic integrals,
\[
 u(a,b)=\frac{(1+r)E(1-\chi)+(1-r)K(1-\chi)}\pi.
\]
These follow by differentiating the circle mean of $|a+be^{i\theta}|$;
at $a=b$ the limit is $2/\pi$. We set $u(0,0)=0$. For conventions and derivative
identities see \url{https://dlmf.nist.gov/19.2} and
\url{https://dlmf.nist.gov/19.4}; DLMF uses the modulus whereas SciPy uses its
square as the parameter.

Subtracting any fixed $\alpha\sqrt T$ from the visible integrand and adding
the exact projection means is unbiased. The means are
$\alpha\sqrt\pi/2$, $-\alpha\sqrt\pi/(4\sqrt2)$, $0$, $0$.
The analogous remote anchor has mean
$\ph(c_4)m\sqrt\pi/(2\sqrt D)$.
These are integration controls, not changes to the response or its noise law.
The four remaining radial integrals are currently evaluated by scrambled Sobol
sampling. Independent scrambles are used in opposite projection factors before
forming bilinear coefficients. Neither a scramble difference nor a stable
optimizer payoff is a deterministic enclosure.

\section*{4. What the computations establish and leave open}
With the new head estimates and the same original-noise kernels, the
twelve-column midpoint common-tail screen is approximately $0.71160588$.
Under assumed head errors $10^{-6}$ and kernel errors $10^{-5}$, the rationalized
price gives a conditional upper screen of about $0.71190294$. Those numerical
response error assumptions remain unproved. The two analytic baseline columns
are independently enclosed by positive rational series. A standard-library
replay checks every front price, the infinite geometric tail, the interval
algebra, and rejection of corrupted data. It proves no pool exclusion.

The new separator was attacked on the droplet in the eleven-dimensional
weighted-degree-$\le7$, charge-one image space of $(Z,G)$ with noise grades
$(1,3)$. Independent evaluations gave payoffs near $0.7742$--$0.7745$. They
are floating evidence against using this separator as a global lower witness.
Their first mixed coefficient is near $0.4813$, so a large adversarial payoff
does not by itself yield a competitive upper score. All four starting points
converged to essentially the same radial visible response; this run supplies no
evidence of a new optimal rounding mechanism.

The rigorous strategic conclusion is narrower and useful: further optimization
within linear Gaussian-phase preprocessors of grades one and three, even with
arbitrary complex phases and convex mixing, cannot beat the incumbent. A viable
upper improvement must leave that class. New dual prices must be attacked in a
larger image space, and the signed full tail must be paid for before reporting
an improvement. The global actual-response inequality remains open.
\end{document}
